\section{Traversal of quasi-subsequential finite-state transducers}

\Cref{alg:qsfst_traversal} illustrates how to deterministically translate an input
word using a $k$-separated quasi-subsequential finite-state transducer.

\begin{algorithm}[h]
\DontPrintSemicolon
\KwIn{$\langle \Sigma, \Omega, Q, X, s, F, \tau, \rho, \eta, \eta_x, \lambda, \psi
	\rangle, \sigma_1 \sigma_2 \ldots \sigma_n$}

$o \gets \epsilon$\;
$q \gets s$\;
\For {$i \in \{ 1, 2, \ldots, n \}$} {
  \For {$j \in \{ 1, 2, \ldots, k \}$} {
    \lIf {$\tau_j(q_j, \sigma_i)$ is \textbf{undefined}} {
      \Return \textbf{undefined}
    }
    $o_j \gets \lambda_j(q_j, \sigma_j)$\;
    $t_j \gets \tau_j(q_j, \sigma_i)$\;
    $r_j \gets \rho_j(t_j)$\;
  }
  $o \gets o \cdot lcp(o_1, o_2, \ldots, o_k)$\;
  $x \gets r_1$\;
  \For {$j \in \{ 2, 3, \ldots, k \}$} {
    $x \gets \eta_x(x, r_j)$\;
  }
  \For {$j \in \{ 1, 2, \ldots, k \}$} {
    $q_j \gets \eta_j(t_j, x)$\;
  }
}
\For {$i \in \{ 1, 2, \ldots, k \}$} {
  \If {$q_i \in F_i$} {
	  \Return $o \cdot \psi_i(q_i)$
  }
}
\Return \textbf{undefined}

\caption{Translation of an input word using a $k$-separated quasi-subsequential
	finite-state transducer.}
\label{alg:qsfst_traversal}
\end{algorithm}

\begin{proposition}
	Let $\mathcal{T} = \langle \Sigma, \Omega, Q, X, s, F, \tau, \rho, \eta,
	\eta_x, \lambda, \psi \rangle$ be a $k$-separa\-ted quasi-subsequential
	finite-state transducer. Let $u = \sigma_1 \sigma_2 \ldots \sigma_n$ be a
	word of length $n \in \mathbb{N}$ over $\Sigma$ and $v$ be a word of length
	$m \in \mathbb{N}$ over $\Omega$ such that $\mathcal{F}_\mathcal{T}(u) = v$.
	Then \Cref{alg:qsfst_traversal} completes on input $\mathcal{T}$ and $u$ in
	$\mathcal{O}(k(n+m))$ time.
\end{proposition}

\begin{proof}
	Lines $1-2$ are executed once and require $\mathcal{O}(1)$ time.
	Lines $5-8$ are executed at most $kn$ times require $\mathcal{O}(1)$ time.
	Line $9$ is executed at most $n$ times and assuming $lcp$ requires
	$\mathcal{O}(ko)$ time where $o = \min_{i=1}^k |o_i|$, the execution over all
	iteration requires $\mathcal{O}(km+n)$ time in total.
	Line $10$ is executed at most $n$ times and requires $\mathcal{O}(1)$ time.
	Line $12$ is executed at most $(k-1)n$ times and requires $\mathcal{O}(1)$
	time.
	Line $14$ is executed at most $kn$ times and requires $\mathcal{O}(1)$ time.
	Line $16$ is executed at most $k$ times and requires $\mathcal{O}(1)$ time.
	Line $17$ is executed at most once and requires $\mathcal{O}(1)$ time.
	Therefore, the total runtime complexity bound for \Cref{alg:qsfst_traversal}
	is $\mathcal{O}(k(n+m))$.
\end{proof}

The complexity bound for \Cref{alg:qsfst_traversal} depends on the ability to compute
longest common prefixes of $k$ words in $\mathcal{O}(km)$ time where $m$ is the
length of the smallest of these words. \Cref{alg:lcp} presents one possible way of
doing that.

\begin{algorithm}[h]
\DontPrintSemicolon
\KwIn{$w_1, w_2, \ldots, w_k$}

$n \gets |w_1|$\;
\For {$j \in \{ 2, 3, \ldots, k \}$} {
	\lIf {$|w_j| < n$} {
		$n \gets |w_j|$
	}
}

$p \gets \epsilon$\;
$i \gets 0$\;

\While {$i \leq n$} {
	\For {$j \in \{ 2, 3, \ldots, k \}$} {
		\lIf {$w_j[i] \neq w_{j-1}[i]$} {
			$\textbf{break}$
		}
	}
	$p \gets p \cdot w_1[i]$\;
	$i \gets i + 1$\;
}

\Return $p$

\caption{Longest common prefix of $k$ words.}
\label{alg:lcp}
\end{algorithm}

\begin{proposition}
	Let $w_1, w_2, \ldots, w_k$ be words over an alphabet $\Sigma$ for $k \in
	\mathbb{N}$ such that $w_i = \sigma_1 \sigma_2 \ldots \sigma_{n_i}$ for $i
	\in \{ 1, 2, \ldots, k \}$. Then \Cref{alg:lcp} completes on input $w_1, w_2,
	\ldots, w_k$ in $\mathcal{O}(kn)$ time where $n = \min_{i=1}^k n_i$.
\end{proposition}

\begin{proof}
	Line $1$ is executed once and requires $\mathcal{O}(1)$ time.
	Line $3$ is executed $k-1$ times and requires $\mathcal{O}(1)$ time.
	Lines $4-5$ are executed once and require $\mathcal{O}(1)$ time.
	Line $8$ is executed at most $(k-1)n$ times and requires $\mathcal{O}(1)$
	time.
	Lines $9-10$ are executed at most $n$ times and assuming that concatenation
	of words requires $\mathcal{O}(1)$ time, then each execution of these lines
	requires at most $\mathcal{O}(1)$ time.
	Therefore, the total runtime complexity bound for \Cref{alg:lcp} is
	$\mathcal{O}(kn)$.
\end{proof}
